\documentclass[11pt]{article}

\usepackage[margin=1in]{geometry}
\usepackage[T1]{fontenc}
\usepackage{lmodern}
\usepackage{booktabs}
\usepackage{longtable}
\usepackage{array}
\usepackage{xcolor}
\usepackage{hyperref}
\usepackage{listings}
\usepackage{enumitem}

\hypersetup{
  colorlinks=true,
  linkcolor=blue,
  urlcolor=blue
}

\definecolor{codebg}{RGB}{246,248,250}
\definecolor{codeframe}{RGB}{210,215,220}

\lstdefinestyle{bridge}{
  basicstyle=\ttfamily\small,
  backgroundcolor=\color{codebg},
  frame=single,
  rulecolor=\color{codeframe},
  breaklines=true,
  columns=fullflexible
}

\title{ProbeBridge\\Technical Manual}
\author{Firmware Project Documentation}
\date{May 2026}

\begin{document}
\maketitle

\begin{abstract}
The ProbeBridge is a USB CDC command adapter for I2C, SPI, UART, and bit-banged 1-Wire targets. It is intended for bench debug, production test fixtures, hardware bring-up, and firmware validation tasks where a small scriptable bus adapter is more useful than a large dedicated instrument. The current implementation targets Raspberry Pi Pico / RP2040 using the Pico SDK.
\end{abstract}

\tableofcontents
\newpage

\section{Project Overview}

The bridge exposes a virtual serial port over USB. A host PC sends newline-terminated text commands such as \texttt{i2c scan}, \texttt{spi xfer 0x9F 0 0 0}, or \texttt{ow reset}. The firmware parses each command, executes the requested bus transaction, and returns a concise status line.

\subsection{Goals}

\begin{itemize}[nosep]
  \item Provide one compact adapter for common low-speed instrument and board interfaces.
  \item Keep the command format readable from both terminals and scripts.
  \item Separate command parsing from hardware backends so parsing can be host-tested.
  \item Make protocol behavior easy to extend without rewriting the console layer.
  \item Demonstrate clean embedded C structure suitable for internal test tooling.
\end{itemize}

\subsection{Non-Goals}

\begin{itemize}[nosep]
  \item It is not a high-throughput logic analyzer.
  \item It is not a voltage-safe universal adapter without external level shifting.
  \item It does not currently implement binary framing, target power control, or firmware updates over the command channel.
\end{itemize}

\section{System Architecture}

\begin{lstlisting}[style=bridge]
PC terminal or automation script
        |
        | USB CDC text commands
        v
Pico SDK stdio USB
        |
        v
line_editor.c
        |
        v
command.c
        |
        v
bus.h abstraction
        |
        v
bus_pico.c
        |
        v
I2C / SPI / UART / 1-Wire targets
\end{lstlisting}

\subsection{Firmware Modules}

\begin{longtable}{>{\ttfamily}p{0.28\linewidth}p{0.62\linewidth}}
\toprule
Module & Responsibility \\
\midrule
src/main.c & Initializes USB stdio, initializes the bus backend, and runs the main command loop. \\
src/line\_editor.c & Collects bytes from the USB CDC stream into newline-terminated command lines. \\
src/command.c & Tokenizes input, validates arguments, dispatches commands, and formats replies. \\
include/bridge/bus.h & Defines the protocol backend interface used by the parser and tests. \\
src/bus\_pico.c & Implements I2C, SPI, UART, and 1-Wire operations using RP2040 peripherals and GPIO timing. \\
tests/parser\_tests.c & Provides fake backend functions for host-side command parser regression tests. \\
\bottomrule
\end{longtable}

\section{Hardware Interface}

\subsection{Default Pico Pinout}

\begin{longtable}{lll}
\toprule
Protocol & Signal & Pico GPIO \\
\midrule
I2C0 & SDA & GP4 \\
I2C0 & SCL & GP5 \\
SPI0 & MISO & GP16 \\
SPI0 & CS & GP17 \\
SPI0 & SCK & GP18 \\
SPI0 & MOSI & GP19 \\
UART0 & TX & GP0 \\
UART0 & RX & GP1 \\
1-Wire & DQ & GP22 \\
\bottomrule
\end{longtable}

\subsection{Electrical Requirements}

\begin{itemize}[nosep]
  \item The RP2040 GPIO bank is 3.3 V only.
  \item I2C and 1-Wire require pull-up resistors. Internal pull-ups are enabled for convenience, but external pull-ups are recommended.
  \item The bridge and target must share ground.
  \item Use level shifting for 1.8 V, 2.5 V, or 5 V targets.
  \item Keep SPI connections short at high clock rates.
\end{itemize}

\section{Command Protocol}

\subsection{Lexical Rules}

Commands are ASCII text lines. Arguments are whitespace separated. Lines beginning with no tokens are ignored, and text after a leading \texttt{\#} token is treated as a comment by the tokenizer.

Numeric arguments support decimal, hexadecimal, octal, and binary notation.

\begin{longtable}{ll}
\toprule
Notation & Example \\
\midrule
Decimal & \texttt{26} \\
Hexadecimal & \texttt{0x1A} \\
Octal & \texttt{032} \\
Binary & \texttt{0b11010} \\
\bottomrule
\end{longtable}

\subsection{Reply Prefixes}

\begin{longtable}{ll}
\toprule
Prefix & Meaning \\
\midrule
\texttt{OK} & Command completed successfully. \\
\texttt{ERR} & Command failed or was rejected. \\
\texttt{TIMEOUT} & Command returned partial data before its timeout expired. \\
\bottomrule
\end{longtable}

\section{Command Reference}

\subsection{General}

\begin{lstlisting}[style=bridge]
help
status
\end{lstlisting}

\subsection{I2C}

\begin{lstlisting}[style=bridge]
i2c scan
i2c speed <hz>
i2c read <addr> <len>
i2c write <addr> <bytes...>
i2c wr <addr> <rx_len> <bytes...>
\end{lstlisting}

I2C addresses are 7-bit addresses in the range \texttt{0x08} through \texttt{0x77}. The \texttt{i2c wr} command performs a repeated-start transaction suitable for register reads.

\subsection{SPI}

\begin{lstlisting}[style=bridge]
spi speed <hz>
spi xfer <bytes...>
spi write <bytes...>
\end{lstlisting}

SPI is full duplex. The bridge returns one received byte for every transmitted byte. The current firmware asserts the fixed SPI chip-select pin for the duration of the transfer.

\subsection{UART}

\begin{lstlisting}[style=bridge]
uart speed <baud>
uart write <bytes...>
uart read <len> [timeout_ms]
\end{lstlisting}

The current UART format is 8N1. The optional read timeout defaults to 100 ms.

\subsection{1-Wire}

\begin{lstlisting}[style=bridge]
ow reset
ow write <bytes...>
ow read <len>
\end{lstlisting}

1-Wire bytes are written and read least-significant bit first according to standard-speed 1-Wire timing.

\section{Build and Flash}

\subsection{Firmware Build}

Install the Pico SDK, set \texttt{PICO\_SDK\_PATH}, then run:

\begin{lstlisting}[style=bridge]
cmake -S . -B build -G Ninja
cmake --build build
\end{lstlisting}

Flash \texttt{build/multi\_protocol\_bridge.uf2} to the Pico.

\subsection{Host Parser Tests}

The parser tests do not require the Pico SDK:

\begin{lstlisting}[style=bridge]
cmake -S . -B build-host -G Ninja -DBRIDGE_HOST_TESTS=ON
cmake --build build-host
.\build-host\parser_tests.exe
\end{lstlisting}

\section{Bring-Up Procedure}

\begin{enumerate}[nosep]
  \item Flash the firmware.
  \item Open the USB CDC serial port.
  \item Run \texttt{help}; confirm the command list is printed.
  \item Run \texttt{status}; confirm default bus speeds.
  \item Run \texttt{i2c scan} with no target connected.
  \item Connect a known I2C target and repeat \texttt{i2c scan}.
  \item Validate SPI with a JEDEC flash ID read.
  \item Validate UART with a TX/RX loopback.
  \item Validate 1-Wire with a DS18B20 or equivalent device.
\end{enumerate}

\section{Testing Strategy}

The project uses two levels of testing:

\begin{itemize}[nosep]
  \item Host parser tests verify command grammar, argument validation, and reply formatting.
  \item Bench tests verify electrical wiring and protocol behavior against known devices.
\end{itemize}

Parser tests should be updated before changing command syntax. Hardware docs should be updated before changing pins, timing assumptions, or electrical behavior.

\section{Future Work}

\begin{itemize}[nosep]
  \item Add configurable SPI mode and chip-select selection.
  \item Add GPIO commands for reset lines and fixture control.
  \item Add ADC voltage measurement for target rail checks.
  \item Add binary framing for high-throughput scripted use.
  \item Add a small adapter PCB with pull-ups, level shifting, and ESD protection.
  \item Add an STM32 backend using the same \texttt{bus.h} contract.
\end{itemize}

\end{document}
